\documentclass[10pt,a4paper]{amsart}
\usepackage[margin=19mm]{geometry}
\usepackage{amsmath,amssymb,mathtools}
\usepackage[scr=rsfs,cal=euler]{mathalfa}
\usepackage{xfrac}
\usepackage[unicode,hidelinks,bookmarks=false]{hyperref}
\newcommand{\ie}{i.e.\ }
\newcommand{\cf}{cf.\ }
\newcommand{\resp}{respectively\ }
\newcommand{\Subsection}{\subsection}
\providecommand{\nobreakdash}{\nobreak\hbox{-}}
\setlength{\emergencystretch}{2em}
\multlinegap0pt
\usepackage{amssymb}
\usepackage{enumerate}
\newtheorem{lemma}{Lemma}
\newtheorem{claim}{Claim}
\newcommand\C{\mathcal C}
\title{Strong log-convexity of genus sequences}
\author{Bojan Mohar}
\date{Source: arXiv:2405.10854v2 (CC BY 4.0)}
\begin{document}
\maketitle

\noindent\emph{Source and licence.} Strong log-convexity of genus sequences. Version arXiv:2405.10854v2 (CC BY 4.0), distributed by arXiv under the Creative Commons Attribution 4.0 International licence, http://creativecommons.org/licenses/by/4.0/, as recorded in the arXiv licence metadata for this exact version and shown on the abstract page https://arxiv.org/abs/2405.10854v2. Full licence notice: \texttt{licenses/arxiv-2405.10854-CC-BY-4.0.txt}.\par\smallskip

\noindent\textit{Editorial scope.} Counted unit: the article's claim bound (source0.tex lines 268--273). The article's complete original proof of it is delivered with it (source0.tex lines 275--328, one \texttt{proof} environment of 54 lines). No mathematics is written, repaired or re-derived: the statement and the proof are the article's own lines, in the article's own notation, and no retained line is substituted or re-flowed.  Retained uncounted as the source-local context the proof needs: lines 143--145; lines 158--160; lines 245--250.  One declared adaptation: exactly 2 ancillary strings inside the retained spans are omitted (image-inclusion commands and, where the source repeats a label, the repeated label definitions) and no other line differs from the source; the figure environments, with their placement, captions and labels, are kept, and the package records the omitted strings in fidelity receipts bound to the affected bodies.  No retained line contains a citation, so the wrapper carries no bibliography and no bibliography entry.  The retained lines contain the article's own diagram code, which the wrapper typesets with the standard tikz packages; no image file is delivered and the package declares no asset dependency.  Editorial formatting only, in addition: an AMS \texttt{amsart} preamble that restates the article's own declarations and macros used by the retained lines, namely the environments \texttt{lemma}, \texttt{claim} on separate counters, together with the standard packages those lines need.  The retained environments print in this document as Lemma 1, Claim 1 in order of appearance, whereas the article prints the same items with the numbering of its own full document.  The complete native source of the article as distributed in the arXiv e-print for this exact version is bundled unchanged as \texttt{sources/159318/source0.tex}.
\begin{lemma}\label{lem:non cofacial faces}
   Suppose that $G$ is a 3-connected graph with a polyhedral embedding in some surface and $\C$ is a sparse collection of its facial cycles. Suppose that $G$ has another embedding $\Pi$ in which none of the cycles in $\C$ is facial. Then the genus of the embedding $\Pi$ is at least $|\C|$.
\end{lemma}

\begin{lemma}\label{lem:disjoint nonhomotopic}
  Suppose that $\C$ is a family of pairwise disjoint and pairwise nonhomotopic and noncontractible cycles on a surface of genus $g\ge1$. Then $|\C|\le 3g-2$.
\end{lemma}

  \begin{claim}\label{cl:lower bound}
     Suppose that $1\le r\le 2k$. Let $q:=\lfloor r/2\rfloor$. Then we have:
     $$
        a_{g+r}(G_{g,k}) \ge 2^{2d(r+q)}\binom{k}{q}.
     $$
  \end{claim}

  \begin{claim}\label{cl:upper bound}
     Suppose that $1\le r\le 2k$. Let $q:=\lfloor r/2\rfloor$ and $\nu=\nu(g,k)$. Then we have:
     $$
        a_{g+r}(G_{g,k}) \le \left((5!)^{31} \nu\right)^{g+2k} \cdot \left( 108d \right)^{36k^2} \cdot 2^{2d(q+r)} .
     $$
  \end{claim}

  \begin{proof}
  We split the vertices of $G=G_{g,k}$ into four sets, $V_A,V_B,V_C,V_D$, which are defined as follows. First, $V_D$ contains for each of the $3k$ connectors the six vertices in the two triangles containing the first and the last edge of the connector (respectively) and the vertex adjacent to that edge but outside of the connector; the set $V_C$ contains the remaining $2d-2$ vertices on each of the connectors; the set $V_B$ contains all vertices that are incident to any of the nontriangular faces and are not in $V_C\cup V_D$; lastly, $V_A$ contains all remaining vertices.
  Note that $|V_D| = 18k$, $|V_C| = 6k(d-1)$, $|V_B| = 2k(m-6)$, and $|V_A| = \nu(g,k)-2k(m+3)$.

  Suppose that $\Pi$ is a 2-cell embedding of $G$ whose genus is (at most) $g+r$. We will compare which facial triangles under $\Pi_0$ are no longer facial under $\Pi$. We say that a $\Pi_0$-facial triangle is \emph{unstable} if it is not $\Pi$-facial; otherwise it is \emph{stable}. We say that a vertex $v\in V(G)$ is \emph{unstable} if it is contained in an unstable triangle. In the first part of the proof, we will show that there are not too many unstable vertices in $V_A\cup V_B\cup V_D$. By ``guessing'' which vertices in $V_A\cup V_B\cup V_D$ are unstable, we get a bound on the number of possible embeddings $\Pi$ having the guessed unstable set. On the other hand, the number of unstable vertices in $V_C$ can be much larger. For the number of rotation systems on the connectors we will need more specific homotopy arguments.

  Let us first consider the set $U$ of all unstable vertices in $V_A\cup V_B$. Let $U'$ be the largest set of unstable triangles, each of which containing at least one vertex in $V_A$ and such that any two triangles in $U'$ are at distance at least 2 in $G$. It is easy to see that the only possible cofacial connector of any two triangles in $U'$ is one of the long faces and such a long face is unique for each triangle in $U'$. (Note that the possible exceptions to this property would be the triangles with their vertices in $V_C\cup V_D$. This was actually the reason to define $V_D$.) Consequently, the set $U'$ is sparse, and Lemma \ref{lem:non cofacial faces} shows that
  $$|U'|\le g+r.$$
  By maximality of $U'$, each vertex in $U$ is at distance at most 2 from $U'$. Since $G$ has maximum degree 7 and vertices of degree 7 are at distance at least 4 from each other, it is easy to see that
  $$|U|\le 30|U'|.$$
  Among the 30 or fewer vertices at distance at most 2 from a triangle in $U'$, there is at most one of degree 7 (by property (d)), and thus they give fewer than $6!(5!)^{29}$ different choices for rotations around these vertices. This ``overcount'' also counts possibilities that stable triangles can have two orientations as the faces in $\Pi$.
  This yields that there are less than
  $$
     \binom{\nu(g,k)}{|U'|}\cdot \left(6 \cdot 6! \cdot (5!)^{29}\right)^{|U'|}
  $$
  ways to select rotations around vertices in $V_A\cup V_B$. Similarly, vertices in $V_D$ have at most $(5!)^{18k}$ ways to choose local rotations around them.

\begin{figure}
\begin{center}

\end{center}
\caption{Triangles on a connector (drawn horizontally). The three black vertices belong to $V_D$.}
\label{fig:3}
\end{figure}

  To count the number of ways to change the local rotations around vertices in $V_C$ and still have genus at most $g+r$, we first note that the above bounds and properties (e) and (f) show that the genus of the subgraph at distance $f(g,k)$ from the long faces of $G$ is at least $g$. Thus we can afford to increase the genus at most by $r$ when changing the rotation around the vertices on the connectors.

  Let us now consider one of the connectors. We will use the notation shown on Fig.~\ref{fig:3}, where the connector is drawn horizontally. Starting on the left with triangle $T_1$, we find the first unstable triangle $T_{i_1}$. Suppose that there is another unstable triangle $T_j$ that is homotopic to $T_{i_1}$ under the embedding $\Pi$. (Here the homotopy is free homotopy among some orientation of two triangles.) We let $T_{i_2}$ be the homotopic triangle with largest possible $i_2$, possibly $i_2=i_1$. Since the two homotopic noncontractible triangles bound a cylinder, any triangle $T_j$ with $i_1\le j\le i_2$ is either contractible or homotopic to $T_{i_1}$. Then we look further and find the smallest $i_3\ge i_2+3$ such that $T_{i_3}$ is unstable. That unstable triangle is disjoint from $T_{i_1}$ and $T_{i_2}$. If $i_3$ exists, then we define $i_4$ and $T_{i_4}$ in the same way as we did to obtain $i_2$ after having fixed $T_{i_1}$. Proceeding, we obtain pairs $(i_1,i_2), (i_3,i_4), \dots, (i_{2t-1},i_{2t})$. The triangles $T_{i_1},T_{i_3},\dots, T_{i_{2t-1}}$ are pairwise disjoint and nonhomotopic, and they increase the genus at most by $r$. Thus, by Lemma \ref{lem:disjoint nonhomotopic}, $t\le 3r-2$.

\begin{figure}
\begin{center}

\end{center}
\caption{Degrees of freedom when the homotopy of unstable triangles is changed. The shaded triangles are facial (with the shown or with the opposite orientation on the surface).}
\label{fig:4}
\end{figure}

  The above paragraph shows that we can determine the number of possibilities of different local rotations on one of the connectors by first finding $i_1,i_2,\dots,i_{2t-1},i_{2t}$ (less than $(2d)^{6r-4}$ possibilities), then estimating that each edge between $T_{i_{2j-1}}$ and $T_{i_{2j}}$ has two choices for joining the two paths on the connector, which amounts all together to at most $2^{2d}$ possibilities, and finally, when changing homotopy, considering the possible ways of local rotations at vertices between $T_{i_{2j}}$ and $T_{i_{2j+1}}$. The latter possibility gives at most $6^3$ ways, for local rotations, including for the three vertices on $T_{i_{2j+1}}$. Figure \ref{fig:4} shows different possibilities for the latter count, where $6^3$ is the worst. The $6^3$ ways can be replaced by $3^3$ to compensate with the overcount of $2^{2d}$ stated above.

  Without trying to optimize and use the same bounds on other connectors, we take it from the proof of Claim \ref{cl:lower bound} that at most $q+r \le 3k$ connectors have unstable triangles. There are at most $2^{3k}$ possible subsets of connectors with unstable triangles. This yields that the number of embeddings is at most
  $$
     2^{3k}\Bigl(2^{2d} \bigl(2d\cdot 3^6\bigr)^{6r-4}\Bigr)^{q+r} .
  $$

  All over, we obtain the upper bound of
  \begin{eqnarray*}
     \binom{\nu}{|U'|} \cdot \left(6 \cdot 6! \cdot (5!)^{29}\right)^{|U'|} && \cdot (5!)^{18k} \cdot
     2^{3k} \Bigl(2^{2d} \bigl(2d\cdot 3^3\bigr)^{6r-4}\Bigl)^{q+r} \\
     &\le& \left(36 (5!)^{30} \nu\right)^{g+2k} \cdot \left( 2 (5!)^6 (54d)^{12k-4} \right)^{3k} \cdot 2^{2d(q+r)}  \\
     &<& \left((5!)^{31} \nu\right)^{g+2k} \cdot \left(108d\right)^{36k^2}
     \cdot 2^{2d(q+r)} ,
  \end{eqnarray*}
  where $\nu=\nu(g,k)$.
  \end{proof}

\end{document}
